%% $Id: gentium-otf-doc.tex 1059 2025-01-08 08:06:10Z herbert $
\listfiles
\documentclass[ngerman,spanish,polish,english,a4paper,paper=a4,DIV=14]{scrartcl}
\usepackage[
  sans=FiraSans-Regular.otf,
  sansFeatures={Scale=MatchUppercase},
  mono=DejaVu Sans Mono,
  monoFeatures={Scale=MatchLowercase,FakeStretch=0.88},
  math=STIXTwo Math,
  mathFeatures={math-style=TeX}]{gentium-otf}
\newfontface\NotoMono[Scale=MatchLowercase,FakeStretch=0.88]{NotoSansMono-Regular.ttf }
\usepackage{babel}
%\usepackage[rmargin=1cm]{geometry}
\usepackage{booktabs}
\usepackage{xltabular}
\usepackage{listings}
\usepackage{showexpl}
\usepackage{xcolor,url,blindtext}
\newcommand\Macro[1]{\texttt{\textbackslash#1}}
\usepackage{hvextern,unicodefonttable}
\usepackage[colorlinks]{hyperref}
\lstset{basicstyle=\ttfamily\small}
\setkeys{hv}{moveToExampleDir,ExampleDir=Examples,showFilename,verbose}

\usepackage[imakeidx]{xindex}
\makeindex[title=Index]

\newenvironment{demoquote}
               {\begingroup
                \setlength{\topsep}{0pt}
                \setlength{\partopsep}{0pt}
                \list{}{\rightmargin\leftmargin}%
                \item\relax}
               {\endlist\endgroup}

\def\Lcs#1{\texttt{\textbackslash#1}\index{#1@\texttt{\textbackslash#1}}}
\def\Largb#1{\texttt{\{}\textit{#1}\texttt{\}}}
\def\Lpack#1{\texttt{#1}\index{#1@\texttt{#1}}\index{Package!#1@\texttt{#1}}}
\def\testfeature#1#2#3{{\fontspec[RawFeature={+#2}]{#1}#3\relax}}
\def\Loption#1{\texttt{#1}\index{#1@\texttt{#1}}\index{Option!#1@\texttt{#1}}}
\def\SUP{\addfontfeatures{RawFeature={+sups}}}

\usepackage{multicol,luacode}
\setlength{\columnsep}{0.3cm}
\setlength{\columnseprule}{0.1pt}

\makeatletter
\blind@addtext{polish}{%
  \def\blindtext@text{%
Honoru myśliwych zaraza Rymsza Libijskich wszystkich Wożny przerywał szczodroty 
opowiadań. Cesarskich znaczy rączy muchom malarstwie spisem różowemi zacnie 
urządzał. Nasze Praga stare żeś Żyje cnoty Tabor. Mówcy pokój Również śmieléj 
wionęła jarzynach liczne drogą jastrząb słowo zabawy odjechał życie. Dano 
ojcu Wyprowadzają kuca dramatycznych myśliwskiém bór pęk żyt maja Wiec. Nię 
Taka Albo tace Usta pęk ucha. }}
\makeatother


\title{Gentium---a complete Greek font with Latin and Cyrillic}
\author{Herbert Voß\thanks{\protect\url{hvoss@tug.org}}}
\usepackage{parskip}
\parindent=0pt


\begin{document}

\maketitle

\tableofcontents


\newpage
\section{The font files}

The \LaTeX\  package \Lpack{gentium-otf} supports the following TrueType fonts:
\begin{verbatim}
GentiumPlus-Regular.ttf            GentiumBookPlus-Regular.ttf  
GentiumPlus-Bold.ttf               GentiumBookPlus-Bold.ttf 
GentiumPlus-Italic.ttf             GentiumBookPlus-Italic.ttf
GentiumPlus-BoldItalic.ttf         GentiumBookPlus-BoldItalic.ttf
\end{verbatim}

The fonts are free available (\url{https://software.sil.org/gentium/})
and also part of any up-to-date distribution, like \TeX Live, MiK\TeX\, or mac\TeX. 

\section{Package options}


Possible optional arguments are

\noindent
\begin{tabular}{@{} >{\ttfamily}l l@{}}
book            &  The book variant of the fonts\\
ScaleRM         & scaling for the serif font, preset to 1\\
defaultfeatures & presetting of features only for rmfamily\\
sans & Sans serif font\\
sansFeatures & Features for the sans serif font\\
mono & Mono font\\
monoFeatures & Features for the mono font\\
math & Math font\\
mathFeatures & Features for the math font\\
\end{tabular}

The Gentium Font comes only with a roman family. If you do not want to load
additional font-packages for a sans-serif, mono and math font, you can define
these fonts with optional arguments. This documentation is typeset with

\begin{verbatim}
\usepackage[
  sans=FiraSans-Regular.otf,
  sansFeatures={Scale=MatchUppercase},
  mono=DejaVu Sans Mono,
  monoFeatures={Scale=MatchLowercase,FakeStretch=0.88},
  math=STIXTwo Math,
  mathFeatures={math-style=TeX}]{gentium-otf}
\end{verbatim}

\newpage

\section{Main font definition}

The main font is defined as:

\begin{verbatim}
\setmainfont{GentiumPlus}[ 
  RawFeature     = {\gentium@figurealign;\gentium@figurestyle},
  Scale          = \gentiumRM@scale ,
  UprightFont    = *-Regular,
  ItalicFont     = *-Italic,
  ItalicFeatures = { SmallCapsFont = *-Italic },
  SlantedFont    = *-Regular,
  SlantedFeatures= {FakeSlant=0.2},
  BoldFont       = *-Bold,
  BoldFeatures   = { SmallCapsFont = *-Bold },
  BoldItalicFont = *-BoldItalic,
  BoldItalicFeatures = { SmallCapsFont = *-BoldItalic },
  BoldSlantedFont= *-Bold,
  BoldSlantedFeatures= {FakeSlant=0.2, SmallCapsFont = *-bold },
  SmallCapsFont  = *-Regular,
%  SmallCapsFeatures={RawFeature=+smcp},
  SmallCapsFeatures={Letters=SmallCaps}, 
  Extension      = .ttf  
]
\end{verbatim}

All available features are listed on \url{https://software.sil.org/gentium/features/}.

\newpage
\section{The default font}
\enlargethispage{3cm}
\selectlanguage{polish}%
\begin{xltabular}{\linewidth}{XX}
\bfseries GentiumPlus & \bfseries GentiumBookPlus\\\hline
Regular\\
\fontspec{GentiumPlus-Regular.ttf}
\blindtext
&
\fontspec{GentiumBookPlus-Regular.ttf}
\blindtext
\\\hline
\strut Bold\\
\fontspec{GentiumPlus-Bold.ttf}
\blindtext
&
\fontspec{GentiumBookPlus-Bold.ttf}
\blindtext
\\\hline
Italic\\
\fontspec{GentiumPlus-Italic.ttf}
\blindtext
&
\fontspec{GentiumBookPlus-Italic.ttf}
\blindtext
\\\hline
BoldItalic\\
\fontspec{GentiumPlus-BoldItalic.ttf}
\blindtext
&
\fontspec{GentiumBookPlus-BoldItalic.ttf}
\blindtext
\end{xltabular}


\newpage
\subsection{Small Caps }
\enlargethispage{3cm}
\begin{tabularx}{\linewidth}{XX}
Default  -- \verb|\scshape| \bfseries GentiumPlus & Book -- \verb|\scshape|\\\hline
Regular\\
\fontspec{GentiumPlus-Regular.ttf}
\scshape\blindtext
&
\fontspec{GentiumBookPlus-Regular.ttf}
\scshape\blindtext
\\\hline
\strut Bold\\
\fontspec{GentiumPlus-Bold.ttf}
\scshape\blindtext
&
\fontspec{GentiumBookPlus-Bold.ttf}
\scshape\blindtext
\\\hline
Italic\\
\fontspec{GentiumPlus-Italic.ttf}
\scshape\blindtext
&
\fontspec{GentiumBookPlus-Italic.ttf}
\scshape\blindtext
\\\hline
BoldItalic\\
\fontspec{GentiumPlus-BoldItalic.ttf}
\scshape\blindtext
&
\fontspec{GentiumBookPlus-BoldItalic.ttf}
\scshape\blindtext
\end{tabularx}


\newpage

\subsection{Slanted -- done by FakeSlant}

\begin{xltabular}{\linewidth}{XX}
\verb|\itshape| \bfseries GentiumPlus &  \verb|\slshape| \bfseries GentiumBookPlus\\\hline
Regular\\
\itshape\blindtext
&
\slshape\blindtext
\\\hline
\strut Bold\\
\itshape\bfseries\blindtext
&
\slshape\bfseries\blindtext
\\\hline
\strut SmallCaps\\
\itshape\scshape\blindtext
&
\slshape\scshape\blindtext
\\\hline
\strut Bold SmallCaps\\
\itshape\scshape\bfseries\blindtext
&
\slshape\scshape\bfseries\blindtext
\\\hline
\end{xltabular}


\selectlanguage{english}%

\newpage



\section{The font features}


\begin{lstlisting}[basicstyle=\ttfamily\small]
bash-3.2$ otfinfo --features GentiumPlus-Regular.otf 
aalt	Access All Alternates
c2sc	Small Capitals From Capitals
ccmp	Glyph Composition/Decomposition
cv13	<unknown feature>
cv14	<unknown feature>
[...]
cv92	<unknown feature>
cv98	<unknown feature>
frac	Fractions
kern	Kerning
liga	Standard Ligatures
mark	Mark Positioning
mkmk	Mark to Mark Positioning
smcp	Small Capitals
ss01	Stylistic Set 1
ss04	Stylistic Set 4
ss05	Stylistic Set 5
ss07	Stylistic Set 7
ss11	Stylistic Set 11
ss12	Stylistic Set 12
subs	Subscript
sups	Superscript
\end{lstlisting}

%$


\subsection{Capitals to Small Caps}

The macro \Lcs{Lctosc}\Largb{arg} is for a local change of \emph{arg} and \Lcs{LCtoSC+} and \Lcs{LCtoSC-} for
a global change of capitals to small caps.

\begin{LTXexample}
Gentium Font ŐŰÉÁÄ 
\Lctosc{Gentium Font ŐŰÉÁÄ}\\
\LCtoSC+ Gentium Font ŐŰÉÁÄ
\end{LTXexample}

\subsection{Capitals to Small Caps and small captitals}

The macro \Lcs{Lctosmcp}\Largb{arg} is for a local change of \emph{arg} and \Lcs{LCtoSMCP+} and \Lcs{LCtoSMCP-}  for
a global change of capitals to small caps.

\begin{LTXexample}
Gentium Font ŐŰÉÁÄ 
\Lctosmcp{Gentium Font ŐŰÉÁÄ}\\
\LCtoSMCP+ Gentium Font ŐŰÉÁÄ 
\end{LTXexample}


\subsection{Ligatures}

The macros \Lcs{Lliga}\Largb{arg} (standard ligatures), \Lcs{Lhlig}\Largb{arg} (historical ligatures), 
\Lcs{Ldlig}\Largb{arg} (discretionary ligatures)
are for a local change of \emph{arg} and \Lcs{LLIGA+}/\Lcs{LLIGA-}, \Lcs{LHLIG+}/\Lcs{LHLIG-}, and \Lcs{LDLIG+}/\Lcs{LDIG-}  for
a global change of capitals to small caps relative to the current group.

\begin{LTXexample}
No ligatures for the upright font!
Schifffahrt ff, fi, ffi, fl

\itshape but for the italic font!
Schifffahrt ff, fi, ffi, fl
{\LLIGA- ff, fi, ffi, fl, Schifffahrt}
\end{LTXexample}

\subsection{Fractions}

There are two macros: \Lcs{Lfrac}\Largb{arg} for a local fraction setting and \Lcs{LFRAC+} and \Lcs{LFRAC-}  for a global
setting relative to the current group.


\begin{LTXexample}[pos=b]
 default: 1/2 1/3 1/4 1/5 \ldots

1/2 \Lfrac{1/3} 1/4 \Lfrac{1/5} 1/6 \Lfrac{1/7} 1/8 \Lfrac{1/9} 1/10 \Lfrac{5/1289}

\LFRAC+
1/2 456/789 1/2 456/789
1/2 1/3 1/4 1/5 1/6 1/7 1/8 1/9 1/10 5/20 9/1289

\LFRAC-
 1/2 1/3 1/4 1/5 1/6 \ldots
\end{LTXexample}



\subsection{Stylistic Sets}

\begin{verbatim}
ss01	Single-story a and g
ss04	Barred-bowl forms
ss05	Slant italic specials
ss07	Low-profile diacritics
ss11	Single-story a (only)
ss12	Single-story g (only)
\end{verbatim}



There is a short command \Lcs{Lssxx}\Largb{text} for the six stylistic sets, 
where xx is the number of the set (two digits)
and \textit{text} the local argument:

\iffalse
\begin{LTXexample}[pos=b]
default: aªàáâãäåāăąǎǟǡǻȁȃȧḁẚạảấầẩẫậắằẳẵặⱥₐᵃ◌ͣgĝğġģǧǵǥḡꞡᵍ\par
ss01: \Lss01{aªàáâãäåāăąǎǟǡǻȁȃȧḁẚạảấầẩẫậắằẳẵặⱥₐᵃ◌ͣgĝğġģǧǵǥḡꞡᵍ}\par
ss11: \Lss11{aªàáâãäåāăąǎǟǡǻȁȃȧḁẚạảấầẩẫậắằẳẵặⱥₐᵃ◌ͣ}\par
ss12: \Lss12{ĝğġģǧǵǥḡꞡᵍ}\par
ss04: đƀǥ \Lss04{đƀǥ}\par
default: aãàáâäåāăǎǟǡǻȁȃȧḁẚảấầẩẫậắằẳẵạặⱥɐæfḟiìíîïĩīĭįǐȉȋḭḯỉịılĺḷḹḻḽꝉₗvṽṿꝟzźżžẑẓẕғӻfiffi\par
ss05: \Lss05{aãàáâäåāăǎǟǡǻȁȃȧḁẚảấầẩẫậắằẳẵạặⱥɐæfḟiìíîïĩīĭįǐȉȋḭḯỉịılĺḷḹḻḽꝉₗvṽṿꝟzźżžẑẓẕғӻfiffi}\par
ss07: áàâǎāãäȧ \Lss05{áàâǎāãäȧ} also for e,o,\ldots
\end{LTXexample}
\fi

\let\myFont\NotoMono
\fvset{fontfamily=myFont}


\begin{externalDocument}[
%  grfOptions={width=0.5\linewidth},
  frame,compiler=lualatex,
  crop,
  force=true,
  usefancyvrb,
  runs=2,code,docType=latex,
  frame,
%  showFilename,
%  align=\centering,
%  lstOptions={columns=flexible}
]{Gentium}
\documentclass{article}
\pagestyle{empty}
\parindent=0pt
%StartVisiblePreamble
\usepackage{gentium-otf}
%StopVisiblePreamble
\begin{document}
default: aªàáâãäåāăąǎǟǡǻȁȃȧḁẚạảấầẩẫậắằẳẵặⱥₐᵃ◌ͣgĝğġģǧǵǥḡꞡᵍ\par
ss01: \Lss01{aªàáâãäåāăąǎǟǡǻȁȃȧḁẚạảấầẩẫậắằẳẵặⱥₐᵃ◌ͣgĝğġģǧǵǥḡꞡᵍ}\par
ss11: \Lss11{aªàáâãäåāăąǎǟǡǻȁȃȧḁẚạảấầẩẫậắằẳẵặⱥₐᵃ◌ͣ}\par
ss12: \Lss12{ĝğġģǧǵǥḡꞡᵍ}\par
ss04: đƀǥ \Lss04{đƀǥ}\par
ss05 is only valid for \verb|\itshape|\\
{\itshape
default: aãàáâäåāăǎǟǡǻȁȃȧḁẚảấầẩẫậắằẳẵạặⱥɐæfḟiìíîïĩīĭįǐȉȋḭḯỉịılĺḷḹḻḽꝉₗvṽṿꝟzźżžẑẓẕғӻfiffi\par
ss05: \Lss05{aãàáâäåāăǎǟǡǻȁȃȧḁẚảấầẩẫậắằẳẵạặⱥɐæfḟiìíîïĩīĭįǐȉȋḭḯỉịılĺḷḹḻḽꝉₗvṽṿꝟzźżžẑẓẕғӻfiffi}}\par
ss07: áàâǎāãäȧ \Lss05{áàâǎāãäȧ} also for e,o,\ldots
\end{document}
\end{externalDocument}

For a global change of the stylistic set one can use the command \Lcs{LSSxx}, where xx
is again the number of the set.

\begin{externalDocument}[
%  grfOptions={width=0.5\linewidth},
  frame,compiler=lualatex,
  crop,
  usefancyvrb,
  force=true,
mpwidth=0.6\linewidth,
  runs=2,code,docType=latex,
  frame,
%  showFilename,
%  align=\centering,
%  lstOptions={columns=flexible}
]{Gentium}
\documentclass{article}
\pagestyle{empty}
%StartVisiblePreamble
\usepackage{gentium-otf}
%StopVisiblePreamble
\begin{document}
đƀǥ $\rightarrow$ \LSS04 đƀǥ
\end{document}
\end{externalDocument}

\subsection{Character variants}
There is a short command \Lcs{Lcvxx}\Largb{text} for the characters, where xx is the number of 
the alternative (two digits)
and \textit{text} the local argument. The macro has an optional argument for passing other values
than 1 to the option (see 43 and 62). Possible variants are:

\begin{verbatim}
cv13	B hook U+0181
cv14	beta U+03B2 U+1D66 U+1D5D
cv17	D hook U+018A
cv19	ezh curl U+0293
cv20	Ezh U+01B7 U+04E0
cv25	rams horn U+0264
cv28	H stroke U+0126
cv37	J stroke hook U+0284
cv43	Eng U+014A
cv44	N left hook U+019D
cv46	Open-O U+0186 U+0254 U+1D10 U+1D53 U+1D97
cv47	OU U+0222 U+0223 U+1D3D U+1D15
cv49	p hook U+01A5
cv55	R tail U+2C64
cv57	T hook U+01AC
cv62	V hook U+01B2 U+028B U+1DB9
cv68	Y hook U+01B3
cv69	Clicks U+01C0 U+01C1 U+01C2 U+2980
cv70	Modifier apostrophe U+02BC U+A78B U+A78C
cv71	Modifier colon U+A789
cv75	Vietnamese-style diacritics U+1EA4...U+1EAB U+1EAE...U+1EB5 U+1EBE U+1EBF U+1EC0...U+1EC5 
                                    U+1ED0...U+1ED7
cv76	Ogonek U+0328 U+0104 U+0105 U+0118 U+0119 U+012E U+012F U+0172 U+0173 U+01EA U+01EB U+01EC U+01ED
cv77	Caron U+010F U+013D U+013E U+0165
cv78	Porsonic circumflex U+0342 U+1F06 U+1F07 U+1F0E U+1F0F U+1F26 U+1F27 U+1F2E U+1F2F U+1F36 U+1F37 
                            U+1F3E U+1F3F U+1F56 U+1F57 U+1F5F U+1F66 U+1F67 U+1F6E U+1F6F U+1F86 U+1F87 
                            U+1F8E U+1F8F U+1F96 U+1F97 U+1F9E U+1F9F U+1FA6 U+1FA7 U+1FAE U+1FAF U+1FB6 
                            U+1FB7 U+1FC0 U+1FC1 U+1FC6 U+1FC7 U+1FCF U+1FD6 U+1FD7 U+1FDF U+1FE6 U+1FE7 
                            U+1FF6 U+1FF7
cv79	Kayan diacritics U+0300 U+0301
cv80	Cyrillic E U+042D U+044D
cv81	Cyrillic shha U+04BB
cv82	Cyrillic breve U+0306
cv83	Capital adscript iota (prosgegrammeni) U+1F88...U+1F8F U+1F98.. U+1F99...U+1F9F U+1FA8...U+1FAF 
                                               U+1FBC U+1FCC U+1FFC
cv84	Serbian and Macedonian Cyrillic alternates U+0431 U+0433 U+0434 U+043F U+0442 U+0453
cv90	Chinantec tones U+02CB U+02C8 U+02C9 U+02CA
cv91	Tone numbers U+02E5 U+02E6 U+02E7 U+02E8 U+02E9 U+A712 U+A713 U+A714 U+A715 U+A716
cv92	Hide tone contour staves U+02E5...U+02E9 U+A712...U+A716
cv98	Empty set U+2205
\end{verbatim}


Don't worry about the tofu in the following listings. Monospace font don't have all
glyphs. The output is okay!


\begin{externalDocument}[
  frame,compiler=lualatex,
  crop,
  force=true,
  usefancyvrb,
  runs=2,code,docType=latex,
  frame,
%  lstOptions={columns=flexible}
]{Gentium}
\documentclass{article}
\pagestyle{empty}
%StartVisiblePreamble
\usepackage[math=STIXTwo Math]{gentium-otf}
%StopVisiblePreamble
\usepackage{parskip}
\begin{document}
13: Ɓ $\rightarrow$\Lcv13{Ɓ}; 14: βᵝᵦ $\rightarrow$\Lcv82{βᵝᵦ}; 17: Ɗ $\rightarrow$\Lcv17{Ɗ}; 
19: ʓ $\rightarrow$\Lcv19{ʓ}; 20: Ʒ $\rightarrow$\Lcv20{Ʒ}; \par
25: ɤ $\rightarrow$\Lcv25{ɤ};  [2]25: ɤ $\rightarrow$\Lcv[2]25{ɤ};
28: Ħ $\rightarrow$\Lcv28{Ħ}; 37: ʄ $\rightarrow$\Lcv37{ʄ}; \par
43: Ŋ $\rightarrow$\Lcv43{Ŋ}; [2]43: Ŋ $\rightarrow$\Lcv[2]43{Ŋ}; [3]43 Ŋ $\rightarrow$\Lcv[3]43{Ŋ};  \par
44: Ɲ $\rightarrow$\Lcv44{Ɲ}; 46: Ɔɔᴐᵓᶗ $\rightarrow$\Lcv46{Ɔɔᴐᵓᶗ}; 47:  Ȣȣᴕᴽ $\rightarrow$\Lcv47{Ȣȣᴕᴽ}; 
49: ƥ $\rightarrow$\Lcv49{ƥ}; 55: Ɽ $\rightarrow$\Lcv55{Ɽ}; 57: Ƭ $\rightarrow$\Lcv57{Ƭ};\par
62: Ʋʋᶹ $\rightarrow$\Lcv62{Ʋʋᶹ}; [2]62: Ʋʋᶹ $\rightarrow$\Lcv[2]62{Ʋʋᶹ};\par
69: ǀǁǂ⦀ $\rightarrow$\Lcv69{ǀǁǂ⦀}; 70: ʼꞋꞌ $\rightarrow$\Lcv70{ʼꞋꞌ}; 71: ꞉ $\rightarrow$\Lcv71{꞉}; \par
75: ẤấẦầẨẩẪẫẮắẰằẲẳẴẵẾếỀềỂểỄễỐốỒồỔổỖỗ $\rightarrow$\Lcv75{ẤấẦầẨẩẪẫẮắẰằẲẳẴẵẾếỀềỂểỄễỐốỒồỔổỖỗ}; \par
76: ĄąĘęĮįŲųǪǫǬǭ $\rightarrow$\Lcv76{Ąą Ęę Įį Ųų Ǫǫ Ǭǭ}; 77: ＾ť $\rightarrow$\Lcv77{＾ť}; \par
78: ἆἇᾆᾇᾶᾷἦἧᾖᾗῆῇἶἷῖῗὖὗῦῧὦὧᾦᾧῶῷἎἏᾎᾏἮἯᾞᾟἾἿὟὮὯᾮᾯ $\rightarrow$\Lcv78{ἆἇᾆᾇᾶᾷἦἧᾖᾗῆῇἶἷῖῗὖὗῦῧὦὧᾦᾧῶῷἎἏᾎᾏἮἯᾞᾟἾἿὟὮὯᾮᾯ};
79: \makebox[0pt]{ò}́ $\rightarrow$\Lcv79{ò́}; \par% U+0300 U+0301
80: Ээ $\rightarrow$\Lcv80{Ээ}; 81: һ $\rightarrow$\Lcv81{һ}; 82: Әә $\rightarrow$\Lcv82{Ә̆ә̆}; \par
83: ᾼᾈᾉᾊᾋᾌᾍᾎᾏῌᾘᾙᾚᾛᾜᾝᾞᾟῼᾨᾩᾪᾫᾬᾭᾮᾯ $\rightarrow$\Lcv83{ᾼᾈᾉᾊᾋᾌᾍᾎᾏῌᾘᾙᾚᾛᾜᾝᾞᾟῼᾨᾩᾪᾫᾬᾭᾮᾯ}; \par
84: бгдптѓ $\rightarrow$\Lcv84{бгдптѓ}; 84 (itshape): {\itshape бгдптѓ $\rightarrow$\Lcv84{бгдптѓ}}; \par
90: ˋˈˉˊ $\rightarrow$\Lcv90{ˋˈˉˊ}; 91: ˥ ˦ ˧ ˨ ˩ ꜒ ꜓ ꜔ ꜕ ꜖ $\rightarrow$\Lcv91{˥ ˦ ˧ ˨ ˩ ꜒ ꜓ ꜔ ꜕ ꜖}; \par
92: ˥ ˦ ˧ ˨ ˩ ꜒ ꜓ ꜔ ꜕ ꜖ (˩˦˥˧˨ ꜖꜓꜒꜔꜕) $\rightarrow$\Lcv92{˥ ˦ ˧ ˨ ˩ ꜒ ꜓ ꜔ ꜕ ꜖ (˩˦˥˧˨ ꜖꜓꜒꜔꜕)}; \par
98: ∅ $\rightarrow$\Lcv98{∅}; 
\end{document}
\end{externalDocument}






Instead of character variation 84 one can also set the script and the language:

\begin{externalDocument}[
  frame,compiler=lualatex,
  crop,
  usefancyvrb,
  force=true,
  runs=2,code,docType=latex,
  frame,
]{Gentium}
\documentclass{article}
\pagestyle{empty}
%StartVisiblePreamble
\usepackage{gentium-otf}
%StopVisiblePreamble
\begin{document}
(upshape): 
{бгдптѓ \addfontfeature{Script=Cyrillic,Language=Serbian}$\rightarrow$ бгдптѓ}

(itshape): \itshape 
бгдптѓ \addfontfeature{Script=Cyrillic,Language=Serbian}$\rightarrow$ бгдптѓ
\end{document}
\end{externalDocument}





\subsection{Numerical alternates}
Only valid for the digits and not for letters!

\begin{LTXexample}
01234\Lsub{a567a}89 01234\Lsup{a567a}89

01234\LSUP+a56789 \LSUP-a01234

01234\LSUB+a56789 \LSUB-a01234
\end{LTXexample}

\newpage

\section{Examples}

\begin{externalDocument}[
  frame,compiler=lualatex,
  crop,
  usefancyvrb,
  force=true,
  runs=2,code,docType=latex,
  frame,
]{Gentium}
\documentclass{article}
\pagestyle{empty}
%StartVisiblePreamble
\usepackage[greek]{babel}
\usepackage{gentium-otf}
\usepackage{blindtext}
%StopVisiblePreamble
\begin{document}
Από τις 140 χώρες που υπάρχουν στο κόσμο, μόνον 19 έχουν πραγματική δημοκρατία. 
Τούτο σημαίνει ότι αυτό το θαυμάσιο και ευαίσθητο πολίτευμα είναι δύσκολο στην 
εφαρμογή του. Και προϋποθέτει ειδικές συνθήκες, χωρίς τις οποίες δεν ημπορεί να 
ανθίσει και να καρποφορήσει.

Δημοκρατία θα ήθελαν να έχουν όλοι οι λαοί. Αλλά δεν την έχουν παρά μόνο εκείνοι 
που, εκτός από την συνειδητοποίηση των προτερημάτων της, σέβονται και τους όρους 
της λειτουργίας της.
\end{document}
\end{externalDocument}

\begin{externalDocument}[
  frame,compiler=lualatex,
  crop,
  usefancyvrb,
  force=true,
  runs=2,code,docType=latex,
  frame,
]{Gentium}
\documentclass{article}
\pagestyle{empty}
%StartVisiblePreamble
\usepackage[russian]{babel}
\usepackage{gentium-otf}
\usepackage{blindtext}
%StopVisiblePreamble
\begin{document}
Из 140 стран мира только 19 имеют реальную демократию. 
Это означает, что эта замечательная и чувствительная система трудно реализовать. 
И это требует особых условий, без которых он не может цветут и принести фрукты.

Демократия хотела бы, чтобы все народы имели. Но у них есть только те, кто, в дополнение 
к реализации его достижений, также уважают условия его работы.
\end{document}
\end{externalDocument}


\newpage

\section{The glyph list}

\displayfonttable{GentiumPlus-Regular.ttf}%[Scale=0.95]

\displayfonttable{GentiumPlus-Italic.ttf}%[Scale=0.95]


\printindex
\end{document}


 